\chapter{Conclusiones y Trabajo Futuro} \label{cap:conclusiones}

\section{Conclusiones}\label{sec:conclusiones}

El objetivo general de este trabajo fue implementar un \textit{pipeline} basado en SLM capaz de identificar errores, notificar y recomendar al desarrollador en el contexto de la DTIC de la Universidad de Cuenca. A continuación se resume el cumplimiento de cada objetivo específico según la evidencia del Capítulo~\ref{cap:resultados}.

\begin{enumerate}
  \item \textbf{Dataset estructurado (objetivo 1).} Se construyó, a partir de 27{,}752{,}836 eventos de \texttt{ms-users}, un conjunto de 223{,}104 eventos relevantes (0.804\,\%) agrupados en 4{,}463 \textit{batches} de 50 eventos, con un \textit{ground truth} por \textit{batch} que contiene 130 errores (2.91\,\%). La etiqueta combina preetiquetado por tres LLM ($\kappa$ de Fleiss $=0.886$), revisión de 329 \textit{batches} por dos expertos, reglas de error derivadas de su criterio y propagación por similitud, con la procedencia de cada etiqueta registrada. La clasificación por tipos de error queda pendiente.
  \item \textbf{SLM adaptados (objetivo 2).} Se implementó y evaluó un conjunto de ocho SLM de distintas familias y tamaños con un \textit{prompt} que comunica el criterio operativo de la DTIC. El ajuste fino LoRA está implementado, pero sus datos de entrenamiento no corresponden a \texttt{ms-users} y el adaptador no fue evaluado contra el \textit{ground truth} final; este objetivo se cumple de forma parcial.
  \item \textbf{Validación (objetivo 3).} Se calcularon \textit{accuracy}, \textit{precision}, \textit{recall}, $F_1$ y AUC de ROC. Contra la referencia por consenso, granite-4.1-3b obtiene el mejor resultado ($F_1=0.951$, AUC $=0.997$). Contra el \textit{ground truth} final corregido por expertos (recálculo propio), todos los modelos tienen \textit{recall} alto (hasta 0.984) pero baja precisión (0.032 a 0.174) y el mejor $F_1$ es el de Qwen3.5-9B (0.295; AUC $=0.987$).
  \item \textbf{Pipeline integrado (objetivo 4).} Se implementó el \textit{pipeline} con prefiltrado, clasificación por \textit{batch}, caché semántico, notificación a Slack, recomendador local y \textit{dashboard}. Falta documentar una prueba con un flujo de logs en vivo de la DTIC.
  \item \textbf{Simulación de errores (objetivo 5).} El mecanismo de inyección está implementado, pero no se dispone todavía de protocolo ni de resultados.
\end{enumerate}

Además de los objetivos, el trabajo deja tres hallazgos:

\begin{itemize}
  \item El acuerdo entre LLM no garantiza corrección: con $\kappa=0.886$, solo el 12.5\,\% de los \textit{batches} de error unánime son errores según los expertos, y el consenso y el \textit{ground truth} final coinciden en el 72.2\,\% de los \textit{batches}.
  \item La clasificación de los SLM es sensible a la referencia: el ordenamiento de los modelos cambia al pasar del consenso al \textit{ground truth} final.
  \item Con $N=20$ y $N=30$, y comparados contra una referencia de 50, el rendimiento de los modelos evaluados cae respecto de $N=50$, aunque en parte por el método de comparación.
\end{itemize}

\section{Limitaciones}\label{sec:limitaciones}

Las principales limitaciones son: revisión experta por una sola persona en cada mitad de la muestra y sin acuerdo entre expertos medido; criterio fijado tras observar los datos; ausencia de una muestra independiente que compruebe las etiquetas propagadas (4{,}134 \textit{batches}); clase positiva pequeña (130); una sola corrida por modelo; ejecución en dos máquinas distintas; y un único microservicio como fuente de datos. Además, las evaluaciones oficiales de la matriz aún usan la referencia por consenso.

\section{Trabajo futuro}\label{sec:trabajo_futuro}

\begin{itemize}
  \item Re-evaluar la matriz completa contra el \textit{ground truth} final con el procedimiento oficial de colecciones y completar las corridas pendientes ($N=20$ de gemma, Qwen3.5 y Ministral; SmolLM2).
  \item Calibrar el umbral sobre $P(\text{error})$ y mejorar el \textit{prompt} para reducir los falsos positivos.
  \item Ajustar con LoRA un SLM con el \textit{ground truth} de \texttt{ms-users} y compararlo con los modelos base.
  \item Medir el acuerdo entre expertos con doble etiquetado de una submuestra y revisar una muestra independiente de lo propagado, siguiendo el esquema de dos equipos y arbitraje de~\cite{sun2026_parsing}.
  \item Probar, como hipótesis a contrastar y no como resultado, tres vías para subir la \textit{precision}, que es el límite principal de los SLM aquí: (i) \textit{few-shot} o RAG con lotes etiquetados por los expertos, separando los ejemplos recuperables del conjunto de evaluación para evitar fuga~\cite{masri2026_slm_severity, zhang2024_logx}; (ii) una cascada con un detector barato y reevaluación con el SLM~\cite{zhang2024_lograg}; y (iii) destilar el criterio experto en reglas obtenidas sin conexión.
  \item Definir y ejecutar el protocolo de inyección de errores y medir la tasa de detección, la latencia de alerta y la calidad de las recomendaciones.
  \item Incorporar métricas y trazas, además de logs, para cubrir la observabilidad completa, y evaluar el esquema colaborativo SLM--LLM para los casos inciertos~\cite{ma2026_adaptivelog}.
  \item Probar el \textit{pipeline} en un flujo en vivo de la DTIC y en otros servicios.
\end{itemize}
